% Options for packages loaded elsewhere
% Options for packages loaded elsewhere
\PassOptionsToPackage{unicode}{hyperref}
\PassOptionsToPackage{hyphens}{url}
\PassOptionsToPackage{dvipsnames,svgnames,x11names}{xcolor}
%
\documentclass[
  12pt]{article}
\usepackage{xcolor}
\usepackage[margin=1in]{geometry}
\usepackage{amsmath,amssymb}
\setcounter{secnumdepth}{5}
\usepackage{iftex}
\ifPDFTeX
  \usepackage[T1]{fontenc}
  \usepackage[utf8]{inputenc}
  \usepackage{textcomp} % provide euro and other symbols
\else % if luatex or xetex
  \usepackage{unicode-math} % this also loads fontspec
  \defaultfontfeatures{Scale=MatchLowercase}
  \defaultfontfeatures[\rmfamily]{Ligatures=TeX,Scale=1}
\fi
\usepackage{lmodern}
\ifPDFTeX\else
  % xetex/luatex font selection
  \setmainfont[]{TeX Gyre Termes}
  \setmathfont[]{texgyretermes-math.otf}
\fi
% Use upquote if available, for straight quotes in verbatim environments
\IfFileExists{upquote.sty}{\usepackage{upquote}}{}
\IfFileExists{microtype.sty}{% use microtype if available
  \usepackage[]{microtype}
  \UseMicrotypeSet[protrusion]{basicmath} % disable protrusion for tt fonts
}{}
\usepackage{setspace}
\makeatletter
\@ifundefined{KOMAClassName}{% if non-KOMA class
  \IfFileExists{parskip.sty}{%
    \usepackage{parskip}
  }{% else
    \setlength{\parindent}{0pt}
    \setlength{\parskip}{6pt plus 2pt minus 1pt}}
}{% if KOMA class
  \KOMAoptions{parskip=half}}
\makeatother
% Make \paragraph and \subparagraph free-standing
\makeatletter
\ifx\paragraph\undefined\else
  \let\oldparagraph\paragraph
  \renewcommand{\paragraph}{
    \@ifstar
      \xxxParagraphStar
      \xxxParagraphNoStar
  }
  \newcommand{\xxxParagraphStar}[1]{\oldparagraph*{#1}\mbox{}}
  \newcommand{\xxxParagraphNoStar}[1]{\oldparagraph{#1}\mbox{}}
\fi
\ifx\subparagraph\undefined\else
  \let\oldsubparagraph\subparagraph
  \renewcommand{\subparagraph}{
    \@ifstar
      \xxxSubParagraphStar
      \xxxSubParagraphNoStar
  }
  \newcommand{\xxxSubParagraphStar}[1]{\oldsubparagraph*{#1}\mbox{}}
  \newcommand{\xxxSubParagraphNoStar}[1]{\oldsubparagraph{#1}\mbox{}}
\fi
\makeatother


\usepackage{longtable,booktabs,array}
\newcounter{none} % for unnumbered tables
\usepackage{calc} % for calculating minipage widths
% Correct order of tables after \paragraph or \subparagraph
\usepackage{etoolbox}
\makeatletter
\patchcmd\longtable{\par}{\if@noskipsec\mbox{}\fi\par}{}{}
\makeatother
% Allow footnotes in longtable head/foot
\IfFileExists{footnotehyper.sty}{\usepackage{footnotehyper}}{\usepackage{footnote}}
\makesavenoteenv{longtable}
\usepackage{graphicx}
\makeatletter
\newsavebox\pandoc@box
\newcommand*\pandocbounded[1]{% scales image to fit in text height/width
  \sbox\pandoc@box{#1}%
  \Gscale@div\@tempa{\textheight}{\dimexpr\ht\pandoc@box+\dp\pandoc@box\relax}%
  \Gscale@div\@tempb{\linewidth}{\wd\pandoc@box}%
  \ifdim\@tempb\p@<\@tempa\p@\let\@tempa\@tempb\fi% select the smaller of both
  \ifdim\@tempa\p@<\p@\scalebox{\@tempa}{\usebox\pandoc@box}%
  \else\usebox{\pandoc@box}%
  \fi%
}
% Set default figure placement to htbp
\def\fps@figure{htbp}
\makeatother


% definitions for citeproc citations
\NewDocumentCommand\citeproctext{}{}
\NewDocumentCommand\citeproc{mm}{%
  \begingroup\def\citeproctext{#2}\cite{#1}\endgroup}
\makeatletter
 % allow citations to break across lines
 \let\@cite@ofmt\@firstofone
 % avoid brackets around text for \cite:
 \def\@biblabel#1{}
 \def\@cite#1#2{{#1\if@tempswa , #2\fi}}
\makeatother
\newlength{\cslhangindent}
\setlength{\cslhangindent}{1.5em}
\newlength{\csllabelwidth}
\setlength{\csllabelwidth}{3em}
\newenvironment{CSLReferences}[2] % #1 hanging-indent, #2 entry-spacing
 {\begin{list}{}{%
  \setlength{\itemindent}{0pt}
  \setlength{\leftmargin}{0pt}
  \setlength{\parsep}{0pt}
  % turn on hanging indent if param 1 is 1
  \ifodd #1
   \setlength{\leftmargin}{\cslhangindent}
   \setlength{\itemindent}{-1\cslhangindent}
  \fi
  % set entry spacing
  \setlength{\itemsep}{#2\baselineskip}}}
 {\end{list}}
\usepackage{calc}
\newcommand{\CSLBlock}[1]{\hfill\break\parbox[t]{\linewidth}{\strut\ignorespaces#1\strut}}
\newcommand{\CSLLeftMargin}[1]{\parbox[t]{\csllabelwidth}{\strut#1\strut}}
\newcommand{\CSLRightInline}[1]{\parbox[t]{\linewidth - \csllabelwidth}{\strut#1\strut}}
\newcommand{\CSLIndent}[1]{\hspace{\cslhangindent}#1}



\setlength{\emergencystretch}{3em} % prevent overfull lines

\providecommand{\tightlist}{%
  \setlength{\itemsep}{0pt}\setlength{\parskip}{0pt}}



 


\makeatletter
\@ifpackageloaded{caption}{}{\usepackage{caption}}
\AtBeginDocument{%
\ifdefined\contentsname
  \renewcommand*\contentsname{Table of contents}
\else
  \newcommand\contentsname{Table of contents}
\fi
\ifdefined\listfigurename
  \renewcommand*\listfigurename{List of Figures}
\else
  \newcommand\listfigurename{List of Figures}
\fi
\ifdefined\listtablename
  \renewcommand*\listtablename{List of Tables}
\else
  \newcommand\listtablename{List of Tables}
\fi
\ifdefined\figurename
  \renewcommand*\figurename{Figure}
\else
  \newcommand\figurename{Figure}
\fi
\ifdefined\tablename
  \renewcommand*\tablename{Table}
\else
  \newcommand\tablename{Table}
\fi
}
\@ifpackageloaded{float}{}{\usepackage{float}}
\floatstyle{ruled}
\@ifundefined{c@chapter}{\newfloat{codelisting}{h}{lop}}{\newfloat{codelisting}{h}{lop}[chapter]}
\floatname{codelisting}{Listing}
\newcommand*\listoflistings{\listof{codelisting}{List of Listings}}
\makeatother
\makeatletter
\makeatother
\makeatletter
\@ifpackageloaded{caption}{}{\usepackage{caption}}
\@ifpackageloaded{subcaption}{}{\usepackage{subcaption}}
\makeatother
\usepackage{bookmark}
\IfFileExists{xurl.sty}{\usepackage{xurl}}{} % add URL line breaks if available
\urlstyle{same}
\hypersetup{
  pdftitle={Tax Evasion and Productivity: Identifying and Correcting for Cost Overreporting},
  pdfauthor={Hans Martinez},
  pdfkeywords={Tax Evasion, Cost Overreporting, Production Function
Estimation, Productivity},
  colorlinks=true,
  linkcolor={black},
  filecolor={Maroon},
  citecolor={black},
  urlcolor={blue},
  pdfcreator={LaTeX via pandoc}}


\title{Tax Evasion and Productivity: Identifying and Correcting for Cost
Overreporting\thanks{I thank Salvador Navarro, David Rivers, and Timothy
Conley for helpful comments. All errors are my own.}}
\author{Hans Martinez\thanks{Department of Economics, Western
University. Email: \href{mailto:hmarti33@uwo.ca}{hmarti33@uwo.ca}. Website: \href{https://hansmartinez.com}{hansmartinez.com}.}}
\date{October 01,
2026\\[0.5em]\href{https://hansmartinez.com/project/tax\_prod/}{\small\textbf{Latest version, click here}}}
\begin{document}
\maketitle
\begin{abstract}
Value-added taxes (VAT) allow firms to deduct the taxes paid on
purchases from the taxes collected on sales, creating an incentive for
firms to overreport the purchases of inputs to reduce their tax bill.
The fundamental issue in detecting input overreporting stems from two
unobservables, true inputs and productivity. Because firms differ in
productivity, a firm reporting high input use, for a given level of
output, may be either overreporting or simply less productive. I use a
structural production-function approach and a benchmark group, a subset
of firms that report inputs correctly, to detect and measure
overreporting. Because the benchmark firms are assumed to share the same
technology as the other firms, I can estimate its parameters. Using
these estimates, I generate predictions and compare how potentially
tax-evading firms would behave without overreporting to what they
report. I apply the method to a Colombian manufacturing survey from 1981
to 1991. During this period, the Colombian sales tax worked as a VAT
and, due to government and market scrutiny, a subset of firms reported
their inputs correctly.

My approach yields a simple yet informative statistical test to detect
input overreporting. Under the null hypothesis of correct reporting, the
test is robust to a wide family of specifications of both the
probability of detection and the cost of overreporting. Among the 5
largest industries, the test rejects the null hypothesis in all 4
industries subject to the sales tax, and, as expected, it fails to
reject in the industry exempt from the sales tax. Across the 7 largest
industries where the null is rejected, the median input overreporting,
as a share of true inputs, amounts to 16 percent. In addition, I find
that, between 1983 and 1987, after a fiscal reform raised the sales tax
in Colombia, firms in affected industries increased the ratio of
overreported to true inputs by about 8.5 percentage points, an implied
median-based overreporting elasticity of the tax rate between 4.7 and
9.5.

The magnitude of detected overreporting matters for methods of
estimating production functions that exploit a flexible input to
identify productivity (Levinsohn and Petrin 2003; Doraszelski and
Jaumandreu 2013; Ackerberg et al. 2015; Gandhi et al. 2020). These
methods typically assume that materials are flexible, and their demand
or first-order conditions reveal productivity. However, materials are
commonly tax deductible and, thus, the input firms have the incentive to
overreport. I show that failing to correct for overreporting yields
output elasticities of materials that are larger, and productivity
distributions that are more dispersed and less persistent. Using the
production-function parameters, I estimate a structural model of firm
input overreporting to ask how does the expected claimed deductions, and
thus government revenue, respond to the tax rate, where expected claimed
deductions include true and undetected fictitious deductions.
Approaching the observed tax rate from the left, the expected deductions
elasticity of the tax rate is -9.7. Undetected overreporting costs the
government around 3 percent of potential sales-tax revenue.
\end{abstract}

\begin{center}
\textbf{JEL Classification:} H26, D22, L60, O47 \\
\textbf{Keywords:} Tax Evasion, Cost Overreporting, Production Function Estimation, Productivity
\end{center}
\bigskip


\setstretch{1.15}
\section*{Introduction}\label{sec-intro}
\addcontentsline{toc}{section}{Introduction}

Firms owned by the highest earners evade more tax through cost
overreporting than firms owned by lower-income earners: relative to the
bottom 80 percent of the income distribution, the tax evaded in Ecuador
by the top 5 percent, as a share of owners' income, is 56 times larger
(Carrillo et al. 2022). Despite its relevance, we still do not know much
about firms' cost overreporting. Tax authorities around the world
document this kind of tax evasion, often through fictitious invoices
(OECD 2017), and the losses they report are significant: between 1.2 and
2.7 percent of total tax revenue in Chile, Colombia, Ecuador, Poland,
and Slovakia, and about a ninth of the entire VAT compliance gap in
Colombia and Slovakia (OECD 2017; Carrillo et al. 2022; CASE and IHS
2018; Longinotti 2022). Because some of these figures are tax
administrations' own estimates, and all are conditional on detected
evasion, true losses are plausibly larger.

Research has documented how individuals' reported income responds to tax
rates, but has overlooked how firms' overreporting of inputs responds.
For individuals, a large literature estimates how reported income
responds to the tax rate and shows that the response includes evasion
and avoidance as well as changes in real activity (Saez et al. 2012;
Chetty 2009). For firms, evidence on responses to tax rates comes mostly
from bunching at kinks and notches in the tax schedule (Best et al.
2015; Almunia and Lopez-Rodriguez 2018) and from misreported imports
(Fisman and Wei 2004). Moreover, economic theory suggests that optimal
commodity-tax rates should be lower on goods whose evasion is more
elastic with respect to the tax rate (Cremer and Gahvari 1993). Yet, to
my knowledge, no study has estimated how firms' overreporting of inputs,
a form of tax evasion that the credit on purchases in value-added taxes
incentivizes, responds to the tax rate. One reason is that true inputs
are unobserved: for a given level of output, a firm reporting high input
use may be overreporting, or it may suffer from low productivity.

I investigate whether input overreporting rises with the tax rate, and
the extent to which it affects government tax revenue through firms'
claimed deductions. Colombia's sales tax worked as a value-added tax,
and in 1983 it raised the rate for most manufacturing industries but
left a few exempt. I find that, between 1983 and 1987, overreporting by
unincorporated firms in the affected industries rose significantly, by
about 8.5 percentage points of true materials. In a structural model
estimated on all industries, I find that a 0.5 percent increase in the
sales-tax rate raises claimed deductions by 16 to 28 percent, while cuts
lower claimed deductions only gradually: the smallest cut
distinguishable from current policy is 8 percent. A back-of-the-envelope
calculation suggests that, as a share of potential net sales-tax revenue
from unincorporated firms in the sample, undetected overreporting costs
about 3 percent. Near the current rate, and holding the tax collected on
sales fixed, the elasticity from the left of this revenue with respect
to the tax rate is about −9 (\textbf{?@sec-app-backofenvelope}).

To distinguish between evasion and productivity, I employ a structural
production-function approach, and then use it to study how evasion
responds to the tax rate. In the absence of evasion, the first-order
condition for materials, the flexible deductible input, reveals the
structural relationship between the firm's materials share of revenue
and the common technology within the industry. The government and market
scrutiny keep corporations from overreporting. Corporations are thus the
benchmark that allows me to identify the common technology in each
industry and the distribution of measurement error in output. Among
unincorporated firms, deviations from the common technology reveal
overreporting up to measurement error in output. Deconvolution then
separates the distribution of overreporting from that of the measurement
error, and the production-function estimates recover productivity net of
overreporting. To study how overreporting responds to the tax rate, I
exploit the 1983 reform, measuring overreporting against the output
elasticity of materials identified from corporations. Finally, employing
an entropy-based method that does not require specifying the
distribution of the unobservables (Schennach 2014), I estimate a
structural model in which firms choose inputs and overreporting to
maximize expected after-tax profits, weighing the tax evaded against the
probability of detection and a cost of evasion that varies with
productivity.

Separating overreporting from productivity also yields three findings of
its own. First, a simple test compares unincorporated firms' materials
shares with the corporate benchmark in each industry: under the null
hypothesis of no overreporting, both should be equal; under the
alternative, overreporting, the difference will be positive. Among the
20 largest industries, the test detects overreporting in 9 of the 18
industries liable for the sales tax, including 4 of the 5 largest, which
together produce about two fifths of manufacturing output. It rejects
decisively, at the 1 percent level, in 7 of them, where, in the median
industry, unincorporated firms overreport an average of 16 percent of
their true materials. It fails to reject in either of the 2 food
industries, which were exempt from the sales tax and thus incentives to
overreport were negligible. The test is valid for a broad class of
models of detection and cost of evasion, and holds under any common
production technology. A rejection is also unambiguous, because no firm
gains by underreporting deductible inputs, so it can only point to
overreporting. Second, failing to correct for overreporting biases
estimates of the production function and of productivity. In those 7
industries, the uncorrected estimates (Gandhi et al. 2020) show output
elasticities of materials that are consistently larger, and productivity
distributions that are more dispersed and less persistent, than the
corrected ones: in the median industry, the ratio of the 90th to the
10th percentile is 66 percent larger and the half-life of productivity
shocks is 52 percent shorter. Third, the structural estimates show a
cost of evasion that is convex in productivity and lowest near the 86th
percentile of productivity, which rationalizes a pattern observed in
other studies: evasion rises with firm size but drastically drops at the
very top (Carrillo et al. 2022).

These results relate to three strands of public economics and to one of
productivity measurement. First, I estimate how a particular form of
evasion by firms responds to the tax rate, a response that the
literature on taxes and revenue has mostly studied through individuals'
reported income. Traditionally, the elasticity of personal taxable
income has combined all the different responses to the rate, evasion
included, into a single statistic (Feldstein 1999; Slemrod and Kopczuk
2002; Saez et al. 2012). However, when part of the cost of evasion is a
transfer rather than a resource cost, this elasticity alone no longer
measures the efficiency cost of taxation (Chetty 2009). Measuring that
cost requires knowing how evasion itself responds to the rate,
separately from real activity. Direct evidence on this response is
scarcer, and comes from the gap between households' consumption and
reported income (Gorodnichenko et al. 2009) and, for firms, from
bunching at kinks and notches in the tax schedule (Best et al. 2015;
Almunia and Lopez-Rodriguez 2018). I provide evidence for firms'
overreporting of inputs under a sales tax that credits purchases, from
both a statutory rate change and a structural model that separates
overreporting from productivity. The structural model also speaks to the
theory of optimal commodity taxation under evasion by estimating how
claimed deductions respond to the rate (Cremer and Gahvari 1993).

Second, I provide measures of firms' cost overreporting without relying
on confidential tax records. Most evidence on firm evasion focuses on
underreported revenue (Slemrod 2019), yet the cost side matters. Firms
offset enforcement on revenue by raising reported costs (Carrillo et al.
2017; Asatryan and Peichl 2017), and paper trails make transactions
between firms harder to misreport than sales to final consumers
(Pomeranz 2015); therefore, firms that sell to other firms have little
room to misreport on the sales side and more on the cost side. Cost
overreporting has mostly been measured directly with transaction-level
tax records in Ecuador and Mexico. These studies reveal the mechanism:
firms acquiring fake invoices from ghost firms (Carrillo et al. 2022;
Zumaya et al. 2021). Tax records, however, are confidential, available
for few countries, and accessible only to selected researchers.
Survey-based measures of firm evasion rely instead on firms'
self-reported share of sales declared to the tax authority (Dabla-Norris
et al. 2019), whose truthfulness cannot be verified. My approach needs
only reported inputs and output, and is thus equally applicable to
standard firm-level surveys, administrative records, and confidential
tax records.

Third, I extend the validation-sample approach from individuals to
firms. This approach infers evasion by comparing an indirect measure of
true income between a truth-reporting group and a group that may
misreport. For example, the self-employed report less income relative to
their food consumption than employees, whose income is reported by their
employers (Pissarides and Weber 1989). Researchers have also used the
same comparison to measure how evasion responds to a tax reform
(Gorodnichenko et al. 2009). For firms, however, the natural indirect
measure, output relative to inputs, also depends on an additional
unobservable, productivity. A firm reporting high input use may be
overreporting or may simply have a negative productivity shock. I use
the production function to separate the two, with corporations as the
truth-reporting group. Unlike the comparisons for individuals, my
approach does not require the two groups to be alike in size or
productivity: corporations are much larger than unincorporated firms,
and the production function accounts for the differences in productivity
behind this gap.

Finally, {the paper} shows how tax-motivated overreporting of inputs
biases production-function estimates. Proxy-variable methods rely on a
flexible input, typically materials, either inverting its demand
(Levinsohn and Petrin 2003; Ackerberg et al. 2015) or using its
first-order condition (Gandhi et al. 2020). These methods treat reported
inputs as true. Work on mismeasured inputs has focused on classical
measurement error (mean zero and independent of the true input), mostly
in capital (Collard-Wexler and De Loecker 2020). Overreporting is
neither: firms choose it based on their own productivity. Furthermore,
it affects exactly the input on which these methods rely. Failing to
correct for overreporting yields output elasticities of materials that
are consistently larger, and productivity distributions that are more
dispersed and less persistent, than the corrected estimates. These
findings matter for any analysis built on measured productivity
differences across firms.

The rest of the paper follows the argument. \textbf{?@sec-model} shows
how the tax rate generates incentives for a firm to overreport
materials, and how its optimal choice depends on the probability of
detection and the cost of evasion. \textbf{?@sec-setting} explains why
Colombian corporations can serve as truth-tellers.
\textbf{?@sec-testing} shows how overreporting can be detected and
measured from production data alone. Then, \textbf{?@sec-pf}
demonstrates how failing to correct for overreporting biases
production-function and productivity estimates. \textbf{?@sec-fiscal}
leverages the 1983 reform to document that overreporting responds to the
tax rate, and \textbf{?@sec-counterfactual} turns the question around
and asks how much revenue the government would lose to overreporting as
it adjusts the tax rate. \textbf{?@sec-conclusion}, as tradition
demands, concludes.

\appendix

\section*{Online Appendix}\label{sec-online-appendix}
\addcontentsline{toc}{section}{Online Appendix}

\protect\phantomsection\label{refs}
\begin{CSLReferences}{1}{1}
\bibitem[\citeproctext]{ref-Ackerberg2015}
Ackerberg, Daniel A, Kevin Caves, and Garth Frazer. 2015.
\emph{IDENTIFICATION PROPERTIES OF RECENT PRODUCTION FUNCTION
ESTIMATORS}. 83: 2411--51. \url{https://doi.org/10.3982/ECTA}.

\bibitem[\citeproctext]{ref-Almunia2018}
Almunia, Miguel, and David Lopez-Rodriguez. 2018. {``Under the Radar:
The Effects of Monitoring Firms on Tax Compliance †.''} \emph{American
Economic Journal: Economic Policy} 10: 1--38.
\url{https://doi.org/10.1257/pol.20160229}.

\bibitem[\citeproctext]{ref-Asatryan2017}
Asatryan, Zareh, and Andreas Peichl. 2017. \emph{Responses of Firms to
Tax, Administrative and Accounting Rules: Evidence from {Armenia}}.
CESifo Working Paper No. 6754. {Center for Economic Studies and ifo
Institute (CESifo)}.

\bibitem[\citeproctext]{ref-Best2015}
Best, Michael Carlos, Anne Brockmeyer, Henrik Jacobsen Kleven, Johannes
Spinnewijn, and Mazhar Waseem. 2015. {``Production Versus Revenue
Efficiency with Limited Tax Capacity: Theory and Evidence from
{Pakistan}.''} \emph{Journal of Political Economy} 123 (6): 1311--55.
\url{https://doi.org/10.1086/683849}.

\bibitem[\citeproctext]{ref-Carrillo2022}
Carrillo, Paul, Dave Donaldson, Dina Pomeranz, and Monica Singhal. 2022.
\emph{Ghosting the Tax Authority: Fake Firms and Tax Fraud}. NBER.
\url{http://www.nber.org/papers/w30242}.

\bibitem[\citeproctext]{ref-Carrillo2017}
Carrillo, Paul, Dina Pomeranz, and Monica Singhal. 2017. {``Dodging the
Taxman: Firm Misreporting and Limits to Tax Enforcement.''}
\emph{American Economic Journal: Applied Economics} 9: 144--64.
\url{https://doi.org/10.1257/app.20140495}.

\bibitem[\citeproctext]{ref-CASE2018}
CASE, and IHS. 2018. \emph{Study and Reports on the {VAT} Gap in the
{EU}-28 Member States: 2018 Final Report}. TAXUD/2015/CC/131. {European
Commission, Directorate General Taxation and Customs Union}.

\bibitem[\citeproctext]{ref-Chetty2009}
Chetty, Raj. 2009. {``Is the Taxable Income Elasticity Sufficient to
Calculate Deadweight Loss? The Implications of Evasion and Avoidance.''}
\emph{American Economic Journal: Economic Policy} 1 (2): 31--52.
\url{https://doi.org/10.1257/pol.1.2.31}.

\bibitem[\citeproctext]{ref-Collard2020}
Collard-Wexler, Allan, and Jan De Loecker. 2020. \emph{Production
Function Estimation and Capital Measurement Error}. NBER.
\url{http://www.nber.org/papers/w22437}.

\bibitem[\citeproctext]{ref-Cremer1993}
Cremer, Helmuth, and Firouz Gahvari. 1993. {``Tax Evasion and Optimal
Commodity Taxation.''} \emph{Journal of Public Economics} 50 (2):
261--75. \url{https://doi.org/10.1016/0047-2727(93)90052-U}.

\bibitem[\citeproctext]{ref-Dabla2019}
Dabla-Norris, Era, Mark Gradstein, Fedor Miryugin, and Florian Misch.
2019. \emph{Productivity and Tax Evasion}. CESifo.
\href{https://www.RePEc.org}{www.RePEc.org}.

\bibitem[\citeproctext]{ref-Doraszelski2013}
Doraszelski, Ulrich, and Jordi Jaumandreu. 2013. {``{R\&D} and
Productivity: Estimating Endogenous Productivity.''} \emph{The Review of
Economic Studies} 80 (4): 1338--83.
\url{https://doi.org/10.1093/restud/rdt011}.

\bibitem[\citeproctext]{ref-Feldstein1999}
Feldstein, Martin. 1999. {``Tax Avoidance and the Deadweight Loss of the
Income Tax.''} \emph{Review of Economics and Statistics} 81 (4):
674--80. \url{https://doi.org/10.1162/003465399558391}.

\bibitem[\citeproctext]{ref-FismanWei2004}
Fisman, Raymond, and Shang-Jin Wei. 2004. {``Tax Rates and Tax Evasion:
Evidence from {`Missing Imports'} in {China}.''} \emph{Journal of
Political Economy} 112 (2): 471--96.
\url{https://doi.org/10.1086/381476}.

\bibitem[\citeproctext]{ref-Gandhi2020}
Gandhi, Amit, Salvador Navarro, and David A Rivers. 2020. \emph{On the
Identification of Gross Output Production Functions}.

\bibitem[\citeproctext]{ref-Gorodnichenko2009}
Gorodnichenko, Yuriy, Jorge Martinez-Vazquez, and Klara Sabirianova
Peter. 2009. {``Myth and Reality of Flat Tax Reform: Micro Estimates of
Tax Evasion Response and Welfare Effects in {Russia}.''} \emph{Journal
of Political Economy} 117 (3): 504--54.
\url{https://doi.org/10.1086/599760}.

\bibitem[\citeproctext]{ref-Levinsohn2003}
Levinsohn, James, and Amil Petrin. 2003. {``Estimating Production
Functions Using Inputs to Control for Unobservables.''} \emph{Review of
Economic Studies} 70 (April): 317--41.
\url{https://doi.org/10.1111/1467-937X.00246}.

\bibitem[\citeproctext]{ref-Pelaez2022}
Longinotti, Fernando Peláez. 2022. \emph{Efficiency and Tax Gap in
{Latin America} and the {Caribbean}: Value Added Tax and Corporate
Income Tax}. CIAT Working Paper No. 1. Inter-American Center of Tax
Administrations (CIAT).

\bibitem[\citeproctext]{ref-OECD2017}
OECD. 2017. \emph{Technology Tools to Tackle Tax Evasion and Tax Fraud}.
OECD. \href{https://www.oecd.org}{www.oecd.org}.

\bibitem[\citeproctext]{ref-Pissarides1989}
Pissarides, Christopher A., and Guglielmo Weber. 1989. {``An
Expenditure-Based Estimate of Britain's Black Economy.''} \emph{Journal
of Public Economics} 39 (June): 17--32.
\url{https://doi.org/10.1016/0047-2727(89)90052-2}.

\bibitem[\citeproctext]{ref-Pomeranz2015}
Pomeranz, Dina. 2015. {``No Taxation Without Information: Deterrence and
Self-Enforcement in the Value Added Tax.''} \emph{American Economic
Review} 105 (8): 2539--69. \url{https://doi.org/10.1257/aer.20130393}.

\bibitem[\citeproctext]{ref-Saez2012}
Saez, Emmanuel, Joel Slemrod, and Seth H. Giertz. 2012. {``The
Elasticity of Taxable Income with Respect to Marginal Tax Rates: A
Critical Review.''} \emph{Journal of Economic Literature} 50 (1): 3--50.
\url{https://doi.org/10.1257/jel.50.1.3}.

\bibitem[\citeproctext]{ref-Schennach2014}
Schennach, Susane M. 2014. {``Entropic Latent Variable Integration via
Simulation.''} \emph{Econometrica} 82: 345--85.
\url{https://doi.org/10.3982/ecta9748}.

\bibitem[\citeproctext]{ref-Slemrod2019}
Slemrod, Joel. 2019. {``Tax Compliance and Enforcement.''} \emph{Journal
of Economic Literature} 57: 904--54.
\url{https://doi.org/10.1257/JEL.20181437}.

\bibitem[\citeproctext]{ref-SlemrodKopczuk2002}
Slemrod, Joel, and Wojciech Kopczuk. 2002. {``The Optimal Elasticity of
Taxable Income.''} \emph{Journal of Public Economics} 84 (1): 91--112.
\url{https://doi.org/10.1016/S0047-2727(01)00095-0}.

\bibitem[\citeproctext]{ref-Zumaya2021}
Zumaya, Martin, Rita Guerrero, Eduardo Islas, et al. 2021.
\emph{Identifying Tax Evasion in Mexico with Tools from Network Science
and Machine Learning}. April. \url{http://arxiv.org/abs/2104.13353}.

\end{CSLReferences}




\end{document}
